\chapter{Appendix 14: Scenario Xi - Algorithmic Architecture}

\section{The Legacy Paradigm \& Its Failures}
Legacy construction utilizes high-entropy, globalized supply chains to build thermodynamically inefficient sprawl. Retrofitting or building alternative structures is blocked by rigid zoning laws, permit bureaucracy, and a lack of localized engineering intelligence, forcing reliance on expensive, centralized legacy firms.

\section{The Incremental Automation Pathway}
\textbf{Phase 1: Manual Trust.} Communities rely on legacy Professional Engineers to manually design and approve localized structures.

\textbf{Phase 2: AI-Assisted Policy.} Parametric algorithms generate CAD models based on local inventory, which are then manually reviewed for thermodynamic and structural integrity.

\textbf{Phase 3: Autonomous Cryptographic Enforcement.} Layer 5 automatically subjects \texttt{ParametricBuildIntent}s to L2 structural and thermal simulations. The Orchestrator refuses to proceed unless the simulation mathematically proves the design is safe and zero-waste, before wrapping the process in a legacy fiat permit shield.

\section{The Human Boundary (Limits of Automation)}
Algorithms can optimize a structure for snow load and passive solar, but they cannot replace the physical "Barn Raising" process. The tactile assembly of timber and 3D-printed joints remains a profound communal human endeavor that cannot be outsourced to robotics.

\section{Orchestration Mechanics (The "How")}
Layer 5 mandates simulation-based policy gates. A \texttt{ParametricBuildIntent} must clear deterministic physics simulations. If successful, it checks for solar rights via Trust Ring consensus. (Triggering Rule 4: Simulation-based Policy Evaluation is missing from the core architecture, detailing how L5 nodes achieve consensus on complex, computationally heavy structural and thermodynamic simulations without relying on a centralized validator).
